\documentclass[11pt,a4paper]{article}
\usepackage[T1]{fontenc}
\usepackage{lmodern}
\usepackage[margin=21mm,headheight=14pt]{geometry}
\usepackage{amsmath,array,tabularx,xcolor,listings,fancyhdr}
\definecolor{Heading}{HTML}{17324A}
\definecolor{Link}{HTML}{18675C}
\definecolor{Muted}{HTML}{465B6A}
\definecolor{Panel}{HTML}{F1F5F7}
\usepackage[colorlinks=true,linkcolor=Heading,urlcolor=Link]{hyperref}
\hypersetup{pdftitle={PSO: Pseudocode and Very Simple Explanations},
  pdfsubject={Base PSO, PSO with DP, and PSO with DP and multithreading}}
\setlength{\parindent}{0pt}
\setlength{\parskip}{5pt}
\setlength{\emergencystretch}{2em}
\renewcommand{\arraystretch}{1.18}
\pagestyle{fancy}
\fancyhf{}
\fancyhead[L]{\small\sffamily\color{Muted}PSO explained simply}
\fancyhead[R]{\small\sffamily\color{Muted}Semester report companion}
\fancyfoot[R]{\small\sffamily\color{Muted}\thepage}
\renewcommand{\headrulewidth}{0.3pt}
\lstset{basicstyle=\ttfamily\fontsize{9.2}{11}\selectfont,
  backgroundcolor=\color{Panel},frame=single,rulecolor=\color{Panel},
  xleftmargin=5pt,xrightmargin=5pt,framesep=5pt,
  columns=fullflexible,keepspaces=true,showstringspaces=false,
  breaklines=true,breakatwhitespace=true,aboveskip=7pt,belowskip=7pt}
\newcommand{\code}[1]{\texttt{#1}}
\newcommand{\simple}[1]{\textbf{In very simple words:} #1}
\raggedbottom

\begin{document}

{\sffamily\bfseries\color{Heading}\LARGE
PSO: pseudocode and\\very simple explanations\par}
\vspace{3mm}
PSO means \emph{particle swarm optimization}. This guide follows the three
programs in \path{Final/}. The pseudocode shows
their main steps in easy language. It is a description to read, rather than
Python code to run.

\section*{1. First, understand the problem}
Imagine that tasks are people and routers are seats. People who talk a lot
should sit close together. Our network also has several floors. A TSV is a
vertical link between floors, like an elevator. Its full name is
\emph{through-silicon via}. We want short routes and we
want to avoid sending most of the traffic through just one elevator.

\begin{tabularx}{\linewidth}{@{}>{\bfseries}p{31mm}X@{}}
Word & Very simple meaning \\
\hline
Particle & One complete possible design: task seats, TSV locations,
TSV choices, and attachment points. \\
Swarm & A group of particles, so we can try several designs. \\
Fitness / score & A number saying how good a design is. Lower is better. \\
Personal best & The best design one particle has found so far. \\
Global best & The best design the whole group has found so far. \\
Velocity & A saved list of swaps to try. Here it is not a physical speed. \\
Iteration & One round of moving and checking all particles. \\
Repair & Fix rules, such as keeping each task exactly once and making TSV
choices valid. \\
\end{tabularx}

\subsection*{What do all three versions try to improve?}
They use the same score. Half comes from communication cost and half from
uneven TSV use, after scaling both numbers to make them comparable.
\[
\text{score}=0.5\,\frac{\text{communication cost}}{\text{cost reference}}
 +0.5\,\frac{\text{TSV imbalance}}{\text{imbalance reference}}.
\]
Communication cost adds \emph{traffic amount $\times$ route cost} for every
communicating pair. If traffic is 4 and the route costs 3, that pair adds 12.
The route uses the configured costs of its links.

TSV imbalance measures how unevenly traffic uses the installed vertical
links. Loads of $[12,0]$ are uneven; $[6,6]$ are even. The code measures upward
and downward traffic separately and uses the larger variance. A link with
zero traffic still counts when calculating this imbalance.

\simple{All three versions search for good designs. DP saves repeated work.
Multithreading gives different remaining checks to different workers.}

\clearpage
\section*{2. Base PSO}
\textbf{Program:} \path{Final/pso_3p5d.py}

\simple{Make a group of designs. Each design learns from its own best attempt
and the group's best attempt. Move, check, remember, and repeat.}

\subsection*{Pseudocode}
\begin{lstlisting}
ALGORITHM BASE_PSO(graph, network, N, maximum_rounds, seed)
    Set the random seed
    Prepare the network, traffic edges, and score references
    particles = MAKE_SWARM(N)              // random, repaired designs
    FOR each particle p, in order
        p.metrics = SCORE_BASE(p.design)   // calculate from scratch
        personal_best[p] = COPY(p)
    global_best = COPY(the lowest-score personal_best)
    stall = 0

    FOR round = 1 TO maximum_rounds
        previous_best_score = global_best.score
        FOR each particle p, in order
            MOVE_AND_REPAIR(p, personal_best[p], global_best)
        FOR each particle p, in order
            p.metrics = SCORE_BASE(p.design)
            IF p.score < personal_best[p].score
                personal_best[p] = COPY(p)
            IF p.score < global_best.score
                global_best = COPY(p)
        IF global_best.score < previous_best_score
            stall = 0
        ELSE
            stall = stall + 1
        IF stall >= 100
            BREAK
    Save and RETURN global_best
\end{lstlisting}

\subsection*{What these steps mean}
\begin{enumerate}
\setlength{\itemsep}{1pt}
\setlength{\parskip}{0pt}
\item \textbf{Make designs.} Start with $N$ different attempts. Repair each one
and start its saved swap lists empty.
\item \textbf{Remember good attempts.} At the start, each particle's first
design is its personal best. The lowest-score one is the global best.
\item \textbf{Move the whole group.} Every particle makes its move using the
best designs known at the start of that round.
\item \textbf{Check and remember again.} Score particles one at a time.
Replace a remembered design only when the new score is strictly lower.
\item \textbf{Stop.} Finish at the round limit, or after 100 rounds with no
new global best. Return the best design found during the run.
\end{enumerate}

For example, if a particle's personal-best score is 10 and its new score is
8, remember the new design. If the group's best score was 9, replace that
too. If its next score is 11, its personal best stays at 8.

\textbf{Why COPY?} A remembered design must stay unchanged when the current
particle moves later. PSO finds the best design it has seen; it does not
prove that no better design exists.

\clearpage
\section*{3. Moving and scoring: shared ideas}
\label{page:shared}
All three versions use the same movement and scoring rules.

\subsection*{A move means trying some swaps}
Suppose a mapping is $[A,B,C,D]$. Swapping positions 0 and 2 gives
$[C,B,A,D]$. Each task still appears once. The real particle also stores
TSV locations, TSV assignments, and attachment points.

\begin{lstlisting}
PROCEDURE MOVE_AND_REPAIR(p, personal_best, global_best)
    Find task swaps toward personal_best and toward global_best
    Try saved old swaps, each with probability 0.10
    Try personal-best task swaps, each with probability 0.04
    Try global-best task swaps, each with probability 0.02
    Save the swaps actually applied as the new task velocity

    For every 3D stack, do the same with its TSV-location swaps
    For each attachment, possibly copy the personal-best choice (0.04)
    Then possibly copy the global-best attachment choice (0.02)

    Repair task placement, TSV locations, and TSV assignments
\end{lstlisting}

Each probability applies to an individual proposed swap or copy. A
probability of 0.10 means about a 10-in-100 chance on each try. It does not
mean that exactly 10 percent of all swaps must happen.

\simple{Keep a little of the previous move, try a little of what worked for
this particle, and try a little of what worked for the group. Then fix any
broken rules.}

\subsection*{How a design gets its score}
The following shows the scoring idea in one combined pass. The actual code
has separate functions for communication cost and TSV loads.

\begin{lstlisting}
FUNCTION SCORE_BASE(design)
    Find the router holding each task
    communication_cost = 0
    Set the load of every installed directed TSV link to zero

    FOR each traffic edge (task_a, task_b, traffic)
        Follow its route using this design's routing rules
        communication_cost += traffic * route_cost
        FOR each directed TSV crossing on that route
            crossing_load += traffic

    up_variance = variance of upward-link loads
    down_variance = variance of downward-link loads
    imbalance = MAX(up_variance, down_variance)
    score = 0.5 * communication_cost / cost_reference
          + 0.5 * imbalance / imbalance_reference
    RETURN (score, communication_cost, imbalance,
            up_variance, down_variance)
\end{lstlisting}

\textbf{Metrics} means the five returned numbers. Later pseudocode writes
\code{p.metrics} when assigning all five, and \code{p.score} when comparing
just the first number. A direction with no TSV links contributes zero
variance.

\clearpage
\section*{4. PSO with DP}
\textbf{Program:} \path{Final/pso_with_only_dp.py}

DP means \emph{dynamic programming}: solve a smaller repeated problem once,
save its answer, and reuse it. This version also remembers whole-design
scores. A \emph{cache} is simply a table of saved answers.

\simple{Use the same PSO search, but check the answer table before doing work
again. The remaining checks run one at a time.}

\subsection*{Pseudocode}
\begin{lstlisting}
ALGORITHM PSO_WITH_DP(graph, network, N, maximum_rounds, seed)
    Set the random seed and prepare the problem
    Read the traffic edges once
    cache = empty whole-design answer table
    particles = MAKE_SWARM(N)
    CACHED_BATCH(particles, cache, workers = NONE)
    personal_best = COPIES of all particles
    global_best = COPY(the lowest-score personal_best)
    stall = 0

    FOR round = 1 TO maximum_rounds
        previous_best_score = global_best.score
        FOR each particle p, in order
            MOVE_AND_REPAIR(p, personal_best[p], global_best)

        CACHED_BATCH(particles, cache, workers = NONE)

        FOR each particle p, in order
            IF p.score < personal_best[p].score
                personal_best[p] = COPY(p)
            IF p.score < global_best.score
                global_best = COPY(p)

        IF global_best.score < previous_best_score
            stall = 0
        ELSE
            stall = stall + 1
        IF stall >= 100
            BREAK

    Save and RETURN global_best
\end{lstlisting}

\subsection*{What changes inside the checking step?}
\code{CACHED\_BATCH}, shown on page~\pageref{page:batch}, first looks up every
design. It groups identical missing designs and calculates each group just
once. With \code{workers = NONE}, those calculations happen one after another.

Inside each missing design's evaluation, DP also remembers reusable route
parts and distances for ordered router pairs. The next page explains the
route part. The score formula remains the one on page~\pageref{page:shared}.

\textbf{Example:} If particles 1 and 3 have exactly the same repaired design,
one evaluation supplies both scores. The same task placement alone is not
enough: TSV locations, TSV assignments, and attachments must match too.
The answer table also identifies the network, traffic, and scoring settings,
because changing any of these can change the score.

\clearpage
\section*{5. The DP part, explained with a tiny route}
Imagine two routes headed toward the same layer:
\begin{lstlisting}
First route:   A -> B -> C -> D
Second route:  X -> B -> C -> D
\end{lstlisting}
After the first route, save the remaining route from B. When the second
route reaches B, reuse that answer. There is no need to follow B, C, and D
again if the routing conditions are the same.

\subsection*{What does the saved answer contain?}
It stores the remaining route cost, the router reached on the target layer,
and the remaining directed TSV crossings. The key says \emph{current router,
target layer, forwarding mode}. Here, the mode says whether to keep the
selected TSV until the next vertical crossing or choose again at each
planar step. The particle's routing settings stay fixed while this table
is used.

\subsection*{Simple route-DP pseudocode}
\begin{lstlisting}
FUNCTION ROUTE_DP(router, target_layer, mode, route_table)
    key = (router, target_layer, mode)
    IF key is in route_table
        RETURN route_table[key]

    IF router is already on target_layer
        answer = (0, router, empty_crossing_list)
    ELSE
        (next_router, step_cost, step_crossings) =
            TAKE_NEXT_ROUTING_STEP(router, target_layer, mode)
        rest = ROUTE_DP(next_router, target_layer, mode, route_table)
        answer.cost = step_cost + rest.cost
        answer.arrival_router = rest.arrival_router
        answer.crossings = step_crossings followed by rest.crossings

    route_table[key] = answer
    RETURN answer
\end{lstlisting}

\simple{The answer from here is the cost of the next step plus the saved
answer from the next place.}

If B to C costs 1 and C to D costs 1, the saved remaining cost at B is 2.
If X to B costs 1, the second route costs $1+2=3$. These are made-up small
costs to show the idea.

The zero at the target layer means the \emph{vertical traversal} is done.
The full route calculation still adds any final movement across that layer
to reach the destination router.

\textbf{Traffic still counts.} If two flows reuse the same saved crossing,
both flows add their own traffic to its load. Saving a route must not make
the second flow disappear from the imbalance calculation.

\textbf{How this matches the program.} This short version calls itself to
make the dependency easy to see. The program walks forward, records the
steps, and fills answers backward instead. It uses fresh route tables for
a design evaluation; the cost and TSV-load passes use separate tables.
The whole-design table lasts across rounds. Thus a changed design does not
accidentally reuse another design's routing answers.

\clearpage
\section*{6. PSO with DP and multithreading}
\textbf{Entry program:} \path{Final/pso_dp_multithreaded.py}\\
It calls the DP program with threaded evaluation enabled.

A \emph{thread} is a worker inside the same running program. The main
thread acts like a teacher handing different papers to assistants to check.
Here, each paper is one complete repaired design.

\simple{First reuse saved answers. Then give different missing designs to
workers. Wait for their scores before starting the next round.}

\subsection*{Pseudocode}
\begin{lstlisting}
ALGORITHM PSO_DP_MULTITHREADED(graph, network, N,
                              maximum_rounds, seed, T)
    Set the random seed and prepare the problem
    Read the traffic edges once
    cache = empty whole-design answer table
    workers = start a pool of T worker threads
    particles = MAKE_SWARM(N)
    CACHED_BATCH(particles, cache, workers)
    personal_best = COPIES of all particles
    global_best = COPY(the lowest-score personal_best)
    stall = 0

    FOR round = 1 TO maximum_rounds
        previous_best_score = global_best.score
        FOR each particle p, in order, on the main thread
            MOVE_AND_REPAIR(p, personal_best[p], global_best)

        CACHED_BATCH(particles, cache, workers)

        FOR each particle p, in order, on the main thread
            IF p.score < personal_best[p].score
                personal_best[p] = COPY(p)
            IF p.score < global_best.score
                global_best = COPY(p)

        IF global_best.score < previous_best_score
            stall = 0
        ELSE
            stall = stall + 1
        IF stall >= 100
            BREAK

    Close the worker pool, including if the run exits with an error
    Save and RETURN global_best
\end{lstlisting}

\textbf{Who does what?} The main thread makes every random move, repairs
designs, manages the answer table, and updates the best designs. Workers
only calculate metrics for different missing designs. Each worker has its
own active particle, TSV assignments, and local route tables.

\textbf{Why wait?} All scores from this round must be ready before the
next round uses the updated global best. Results are collected in submitted
order and best designs are updated in particle order. This keeps the seeded
DP search repeatable across worker counts for the same inputs.

The pool is reused across rounds. With one effective worker, the program
checks designs directly on the main thread. Unless configured otherwise,
the worker count is the smallest of the CPU count, swarm size, and four.

\clearpage
\section*{7. The checking step used by both DP versions}
\label{page:batch}
\textbf{Cache hit} means ``we already have the answer.'' A \textbf{cache miss}
means ``we must calculate it.'' This pseudocode assumes the whole-design
cache is enabled, as it is by default.

\begin{lstlisting}
PROCEDURE CACHED_BATCH(particles, cache, workers)
    pending = empty groups, kept in first-seen order

    FOR each particle p, in order
        key = (problem_settings, complete_repaired_design_of_p)
        IF key is in cache
            p.metrics = cache[key]
        ELSE IF key is already in pending
            Add p to that key's group
        ELSE
            Create a group with COPY(p.design) and particle p

    IF workers is NONE
        Calculate SCORE_DP for each pending group, one at a time
    ELSE
        Send one SCORE_DP job per pending group to the workers
        Wait for all jobs; collect results in submitted order

    FOR each pending group and its returned metrics, in order
        IF cache has room
            Save (group.key, metrics) in cache
        Give those metrics to every particle in the group

FUNCTION SCORE_DP(design)
    Use the same scoring rules as SCORE_BASE
    Use fresh route-DP tables and an ordered router-pair table
    RETURN the five metrics
\end{lstlisting}

The real whole-design cache holds up to 200,000 entries. When full, it
continues reading old answers but stops adding new ones. Duplicate missing
designs within a batch can still share one evaluation.

\subsection*{A small example comparing all three}
Suppose a batch has six particles: $[A,B,A,C,B,C]$. Here A, B, and C mean
three complete designs, and the answer table starts empty.

\begin{tabularx}{\linewidth}{@{}>{\bfseries}p{30mm}p{24mm}X@{}}
Version & Full checks & What happens \\
\hline
Base PSO & 6 & Check each particle separately. \\
PSO with DP & 3 & Check A, B, and C once each, one after another. \\
DP + threads & 3 & Give A, B, and C to workers; reuse each result. \\
\end{tabularx}

With $T$ workers and $U$ different missing designs, at most $\min(T,U)$ jobs
can be in progress together. Three checks do not guarantee three times the
speed: scheduling and serial work take time, and Python's usual interpreter
lock can limit CPU-heavy threads. Measure elapsed time to see the benefit.

\subsection*{Three sentences to remember}
\textbf{Base PSO:} Try designs and learn from the best ones.\\
\textbf{PSO with DP:} Do that while reusing answers to repeated calculations.\\
\textbf{PSO with DP and threads:} Reuse answers, then share the remaining
design checks among workers.

\end{document}
